\chapter{Eyelid Inflammation (Blepharitis)}

\section*{Definition}
Blepharitis is a common, chronic inflammation of the eyelid margins. It affects people of all ages and is frequently associated with dry eye syndrome and meibomian gland dysfunction.

\section*{What Happens in Your Body}
Blepharitis is classified as anterior (affecting the eyelash follicles) or posterior (affecting the meibomian glands). Anterior blepharitis is often caused by staphylococcal colonization or seborrheic dermatitis. Posterior blepharitis involves meibomian gland dysfunction with abnormal lipid secretion, leading to tear film instability and ocular surface inflammation.

\section*{Causes}
\begin{itemize}
\item Staphylococcal bacterial colonization of lid margins
\item Seborrheic dermatitis
\item Meibomian gland dysfunction
\item Rosacea (ocular rosacea)
\item Demodex mite infestation
\item Contact lens wear
\item Allergies
\item Dry eye syndrome
\end{itemize}

\section*{When to See a Doctor}
\begin{itemize}
\item Red, swollen, and irritated eyelid margins
\item Crusting and flaking of the eyelids (especially on waking)
\item Burning or gritty sensation in the eyes
\item Itching of the eyelids
\item Watery eyes
\item Crusty eyelashes or loss of eyelashes
\item Sensitivity to light
\item Chalazion or hordeolum (stye) formation
\end{itemize}

\section*{Related Symptoms}
\begin{itemize}
\item Dry eyes (dry eye syndrome)
\item Conjunctivitis (pink eye)
\item Stye (hordeolum)
\item Chalazion
\item Eyelash loss (madarosis)
\item Ocular rosacea
\end{itemize}

\section*{Self-Care / Home Management}
\begin{itemize}
\item Apply warm compresses for 5-10 minutes daily
\item Gently scrub the eyelid margins with a clean washcloth or lid scrub
\item Use OTC artificial tears for dry eye symptoms
\item Avoid eye makeup during flare-ups
\item Clean the eyelids with baby shampoo diluted in warm water
\item Replace eye makeup every 3 months
\end{itemize}

\section*{How Your Doctor Will Evaluate You}
\begin{itemize}
\item Slit-lamp examination of the eyelid margins
\item Evaluation of meibomian gland function
\item Eyelash microscopy for Demodex
\item Assessment of tear film break-up time
\item Culture of lid margins if severe or treatment-resistant
\end{itemize}

\section*{Treatment Options}
\begin{itemize}
\item Warm compresses and eyelid hygiene (cornerstone of treatment)
\item Antibiotic ointments (erythromycin, bacitracin) for bacterial blepharitis
\item Topical cyclosporine for ocular rosacea
\item Oral doxycycline or minocycline for meibomian gland dysfunction
\item Topical corticosteroid drops for acute inflammation
\item Management of associated conditions (dry eye, rosacea)
\end{itemize}
\section*{References}
\begin{thebibliography}{99}
\bibitem{eyelid_inflammation:ref1} Lindsley K, Matsumura S, Hatef E, et al. ``Interventions for Chronic Blepharitis.'' Cochrane Database Syst Rev. 2012;(5):CD005556.
\bibitem{eyelid_inflammation:ref2} Putz RA, Putz CE. ``Blepharitis: A Review of Current Treatment Options.'' J Ophthalmic Vis Res. 2014;9(4):480--483.
\bibitem{eyelid_inflammation:ref3} Bernardes TF, Bonfioli AA. ``Blepharitis.'' Semin Ophthalmol. 2010;25(3):79--83.
\bibitem{eyelid_inflammation:ref4} Jackson WB. ``Blepharitis: Current Strategies for Diagnosis and Management.'' Can J Ophthalmol. 2008;43(2):170--174.
\bibitem{eyelid_inflammation:ref5} Kanski JJ, Bowling B. ``Clinical Ophthalmology: A Systematic Approach.'' 7th ed. Elsevier; 2011.
\end{thebibliography}
